\newcommand{\DoNotLoadEpstopdf}{}
\documentclass[11pt]{article}
\usepackage[T1]{fontenc}
\usepackage{lmodern}
\usepackage[margin=1in]{geometry}
\usepackage{amsmath,amssymb,amsthm,mathtools}
\usepackage{microtype}
\usepackage{enumitem}
\usepackage{booktabs,array}
\usepackage{xcolor}
\usepackage{tikz}
\usetikzlibrary{arrows.meta,positioning,calc,decorations.pathreplacing}
\usepackage{needspace}
\usepackage[hidelinks,unicode]{hyperref}
\input{glyphtounicode}
\input{glyphtounicode-cmex}
\pdfgentounicode=1
\pdftrailerid{}
\hypersetup{pdftitle={Random Walks on Critical FK-Ising Maps and Liouville Brownian Motion},pdfauthor={OpenAI}}
\newtheorem{theorem}{Theorem}[section]
\newtheorem{proposition}[theorem]{Proposition}
\newtheorem{lemma}[theorem]{Lemma}
\newtheorem{corollary}[theorem]{Corollary}
\theoremstyle{definition}
\newtheorem{definition}[theorem]{Definition}
\theoremstyle{remark}
\newtheorem{remark}[theorem]{Remark}
\newcommand{\R}{\mathbb R}
\newcommand{\N}{\mathbb N}
\newcommand{\C}{\mathbb C}
\newcommand{\Sph}{S^2}
\newcommand{\PP}{\mathbb P}
\newcommand{\EE}{\mathbb E}
\newcommand{\E}{\mathcal E}
\newcommand{\dd}{\,\mathrm d}
\newcommand{\HS}{\mathrm{HS}}
\newcommand{\ind}{\mathbf 1}
\newcommand{\norm}[1]{\lVert#1\rVert}
\newcommand{\abs}[1]{\lvert#1\rvert}
\newcommand{\ip}[2]{\langle#1,#2\rangle}
\DeclareMathOperator{\diam}{diam}
\DeclareMathOperator{\dist}{dist}
\DeclareMathOperator{\supp}{supp}
\DeclareMathOperator{\Law}{Law}
\DeclareMathOperator{\Id}{Id}
\setlist[enumerate]{itemsep=4pt,topsep=6pt}
\setlist[itemize]{itemsep=4pt,topsep=6pt}
\numberwithin{equation}{section}
\title{Random Walks on Critical FK--Ising Maps\\and Liouville Brownian Motion}
\author{OpenAI}
\date{October 5, 2026}
\begin{document}
\maketitle
\begin{abstract}
We prove that stationary random walks on critical spherical FK--Ising
planar maps converge to Liouville Brownian motion on the ordinary unit-area
$\sqrt3$-quantum sphere, using the geometric and electrical results of the
spectral companion. The walk chooses uniformly among all incident
half-edges, retaining loops and multiple edges. With the corner measure as
the stationary law, the deterministic time acceleration is exactly the
number of map edges. For any deterministic metric scale giving the
metric-measure limit, convergence retains that same surface and any fixed
finite number of conditionally independent walks with their time parameters.
\end{abstract}
\section{Introduction}\label{sec:introduction}

The scaling limit of a random surface does not by itself determine the
scaling limit of a walk on that surface. A time-parametrized walk also
depends on the electrical energy and on the mass assigned to vertices.
We study this additional limit for critical Fortuin--Kasteleyn (FK)--Ising
maps, using the common-coordinate and electrical convergence results of
the spectral companion~\cite{SC}.

\subsection{The model and the result}

Let $\mathcal F_n$ be the finite set of rooted pairs $(M,A)$, where $M$
is a connected planar multigraph with $n$ edges embedded in the oriented
sphere and $A\subseteq E(M)$. Rooting is at an oriented edge, and pairs
are identified by orientation-preserving homeomorphisms preserving the
root and $A$. Loops and multiple edges are allowed. If $M^*$ is the dual
map, set
\[
 A^*=\{e^*:e\notin A\},\qquad
 \ell(M,A)=k(A)+k(A^*)-1,
\]
where each component count includes isolated vertices of the corresponding
spanning subgraph. The critical spherical FK--Ising law is
\begin{equation}\label{eq:fk-law}
 \PP\bigl((M_n,A_n)=(M,A)\bigr)
 =\frac{2^{\ell(M,A)/2}}{Z_n},\qquad
 Z_n=\sum_{(M,A)\in\mathcal F_n}2^{\ell(M,A)/2}.
\end{equation}
Write $V_n=V(M_n)$, let $d_n$ be graph distance with every original edge
of length one, and let
\begin{equation}\label{eq:corner-measure}
 \mu_n(v)=\frac{\deg_{M_n}(v)}{2n}.
\end{equation}
A loop contributes two to the degree. Conditional on $(M_n,A_n)$, a
simple random walk chooses an incident half-edge uniformly and traverses
its edge. Its initial vertex has law $\mu_n$, which is stationary.
The walk uses all map edges, not just those in $A_n$.

Let $(S,h)$ be the ordinary unit-area $\sqrt3$-Liouville quantum gravity
(LQG) sphere: the quantum-sphere law disintegrated at area one, with its
distinguished points forgotten. Write $\mu_h$ for its area probability
measure and $D_h$ for its intrinsic metric in a fixed deterministic
normalization. Liouville Brownian motion is the diffusion associated with
the closure on $L^2(\mu_h)$ of
\begin{equation}\label{eq:continuum-form}
 \E_h(f,f)=\frac12\int_S \abs{\nabla_g f}^2\dd\mathrm{vol}_g,
 \qquad f\in C^\infty(S),
\end{equation}
where $g$ is any smooth metric in the conformal class of the sphere.
Conformal invariance of the two-dimensional Dirichlet integral makes
this definition independent of $g$. The factor $1/2$ fixes the time
convention.

For a walk $X^n$, let $\overline X^n$ be its continuous interpolation on
the cable realization $\overline M_n$, in which each edge is a unit
interval and each selected edge is traversed linearly in one unit of
time. We also write $d_n$ for the cable path metric. We use the
curve-decorated Gromov--Hausdorff--Prokhorov topology: after isometric
embeddings into common metric spaces, the underlying spaces converge in
Hausdorff distance, the measures in Prokhorov distance, and every curve
uniformly on each bounded time interval.

\begin{theorem}\label{thm:main}
Let $(M_n,A_n)$ have law \eqref{eq:fk-law}, and let $a_n>0$ be any
deterministic sequence with $a_n\to0$ such that
\begin{equation}\label{eq:assumed-ghp}
 (V_n,a_nd_n,\mu_n)\ \Longrightarrow\ (S,D_h,\mu_h)
\end{equation}
in Gromov--Hausdorff--Prokhorov law. Given the map, start $r$ independent
simple random walks $X^{n,1},\ldots,X^{n,r}$ independently from $\mu_n$.
Given $h$, start $r$ independent Liouville Brownian motions
$B^{h,1},\ldots,B^{h,r}$ independently from $\mu_h$, with the convention
\eqref{eq:continuum-form}. For every fixed integer $r\ge1$,
\begin{equation}\label{eq:main-limit}
 \left(\overline M_n,a_nd_n,\mu_n;
       (\overline X^{n,j}_{nt})_{t\ge0,\,1\le j\le r}\right)
 \ \Longrightarrow\
 \left(S,D_h,\mu_h;
       (B^{h,j}_t)_{t\ge0,\,1\le j\le r}\right)
\end{equation}
in the curve-decorated Gromov--Hausdorff--Prokhorov topology.
\end{theorem}

Thus the deterministic acceleration can be taken to be exactly $n$.
The finite-product assertion retains the conditional walk law in the
random environment, rather than only the distribution of one trajectory
after averaging over the surface.

The proof uses more than the metric-measure premise
\eqref{eq:assumed-ghp}. The spectral companion supplies a joint conformal
coordinate, unit conductivity, uniform local energy recovery, and
convergence of centered inverse operators. We state these inputs
precisely in Proposition~\ref{prop:inputs}; their proofs and the
foundational companion results on which they depend are not reproduced
here. The path-convergence argument from those inputs is proved in full.

\subsection{History and related work}

Sheffield's inventory-accumulation encoding related critical FK-decorated
planar maps to correlated Brownian paths~\cite{Sheffield2016Inventory}.
Gwynne and Sun proved that the finite-volume FK contour encoding converges
to correlated Brownian motion conditioned to stay in the quadrant and
return to the origin~\cite[Theorem~1.8]{GwynneSun2015Finite}.
The mating-of-trees theory of Duplantier, Miller, and Sheffield, and its
finite-volume sphere form due to Miller and Sheffield, identify the
corresponding continuum quantum
surfaces~\cite{DuplantierMillerSheffield2021Mating,MillerSheffield2019SphereMating}.
The Brownian paths in that encoding describe boundary lengths; the vertex
walk considered here is a different stochastic process.
Liouville Brownian motion was constructed by Garban, Rhodes, and Vargas
and by Berestycki~\cite{GarbanRhodesVargas2016LBM,Berestycki2015Diffusion}.
Its Dirichlet-form description is developed in
\cite[Theorem~1.7]{GarbanRhodesVargas2014Dirichlet}.

For mated-CRT maps, Gwynne, Miller, and Sheffield proved quenched walk
convergence modulo time
parametrization~\cite[Theorem~3.4]{GwynneMillerSheffield2021Tutte}.
Berestycki and Gwynne subsequently retained the
parametrization~\cite[Theorems~1.2 and~1.4]{BerestyckiGwynne2022MatedCRTLBM},
using a deterministic median-exit clock comparable to inverse cell area.
Its asymptotic constant is not identified there
\cite[Remark~1.3]{BerestyckiGwynne2022MatedCRTLBM}.
The setting here is instead the ordinary finite FK--Ising sphere, the
walks start in stationarity, and the limit is joint with the intrinsic
metric and area measure. The exact clock in Theorem~\ref{thm:main} rests
on the unit-conductivity theorem of the spectral
companion~\cite[Theorem~5.2]{SC}.

\subsection{From operators to paths}

The principal issue is tightness with the time parameter retained.
Convergence of eigenvalues, or even convergence of the finite-dimensional
walk laws, does not exclude rapid excursions of appreciable spatial
diameter. The companion's convergence of centered inverses gives
positive-time semigroup convergence. Uniform tracking of the exploration
tour, which lists the graph corners and parametrizes quantum area in the
limit, then identifies the finite-dimensional distributions in one common
coordinate.

To control entire paths, we use smooth spatial cutoffs. Local energy
recovery provides uniformly accurate graph cutoffs with asymptotically
no excess energy. Resolvent smoothing gives test functions whose
generators are uniformly bounded at each fixed smoothing parameter.
A reversible forward--backward martingale decomposition converts the
form-norm error between these two approximations into an error uniform
along a stationary trajectory. Optional sampling of the smoothed
cutoffs then gives a stopping-time tightness estimate.

The local contribution is this passage from the companion's electrical
and operator limits to stationary, time-parametrized path convergence.
The maximal inequality is a finite-state analogue of the classical
Lyons--Zheng forward--backward
decomposition~\cite{LyonsZheng1988}; we prove the jump-chain identity
directly. The last tightness step uses Aldous's
criterion~\cite{Aldous1978}.
Neither an additional mixing assumption nor a heat-kernel estimate is
needed for this passage. Stationarity is essential to the argument as
given; we do not assert convergence from arbitrary starting vertices.

Section~\ref{sec:inputs} records the precise companion inputs and checks
the normalization. Section~\ref{sec:operators} obtains the
finite-dimensional laws. Section~\ref{sec:tightness} proves the stationary
maximal estimate and path tightness. Section~\ref{sec:paths} transfers the
limit to cable paths, removes the exponential clock, and assembles the
joint convergence in random environments.

\section{Energy normalization and the common coupling}\label{sec:inputs}

The proof needs more than the metric-measure limit: it needs convergence of
inverse operators and uniform recovery of spatial test functions with the
correct energy.  We state these inputs from the spectral companion
\cite{SC} precisely, keeping its unhalved electrical convention separate
from the Brownian convention of the theorem.

\subsection{The graph energy and the spatial scale}

Let \(P_n\) be the one-step transition operator of the walk, let
\(L_n=P_n-I\), and put
\[
 G_n=-nL_n=n(I-P_n).
\]
Thus \(-G_n\) is the generator with total attempt rate \(n\), including
attempts along loops.  It is self-adjoint on \(L^2(\mu_n)\).  For a real
vertex function \(u\), write
\(\E_n(u)=\ip{u}{G_nu}_{L^2(\mu_n)}\), and use polarization for the
bilinear form.  Choose either orientation of each edge and write
\(d_eu=u(e^+)-u(e^-)\).

\begin{lemma}[Graph energy]\label{lem:graph-energy}
For all real vertex functions \(u,v\),
\begin{equation}\label{eq:graph-energy}
 \E_n(u,v)=\frac12\sum_{e\in E(M_n)}d_eu\,d_ev.
\end{equation}
Every parallel edge is a separate summand, every loop has zero summand,
and the kernel of \(G_n\) consists of the constants.
\end{lemma}

\begin{proof}
Since \(\mu_n(x)=\deg(x)/(2n)\), every incident half-edge has
\(\mu_n\)-weighted accelerated rate \(1/2\).  A nonloop edge with
endpoints \(x,y\) therefore contributes
\[
 \frac12\bigl(u(x)(v(x)-v(y))+u(y)(v(y)-v(x))\bigr)
 =\frac12d_eu\,d_ev
\]
to \(\ip{u}{G_nv}_{L^2(\mu_n)}\).  Both half-edges of a loop contribute
zero.  Summing proves \eqref{eq:graph-energy}, including multiplicities.
Zero energy forces equality across each edge, hence constancy by
connectedness.
\end{proof}

On the limiting conformal sphere, write \(\E_h(f)=\E_h(f,f)\), also for
the closure, and let \(G\) be its nonnegative self-adjoint operator.  The
existence and identification of this closure are part of
\cite[Proposition~8.3]{SC}.  In the companion's notation the unhalved
operators are
\begin{equation}\label{eq:raw-operators}
 B_n=2G_n=-2nL_n,\qquad B_h=2G.
\end{equation}
In particular, the inverses imported below are inverses of \(2G_n\) and
\(2G\), not of \(G_n\) and \(G\).

The companion fixes space by a deterministic diameter normalization.  The
next elementary comparison, reproduced from \cite[Lemma~2.3]{SC}, shows
why its conclusions apply to the sequence \(a_n\) in the theorem.
For a finite real random variable \(Y\), let \(\mathfrak m(Y)\) be the
unique zero of \(r\mapsto\EE[\arctan(Y-r)]\).  Set
\[
 \Delta_n=\max\{1,\operatorname{diam}(V_n,d_n)\},\qquad
 \Delta=\operatorname{diam}(S,D_h),
\]
and define
\begin{equation}\label{eq:diameter-normalization}
 \widehat a_n=
 \exp\{\mathfrak m(\log\Delta)-\mathfrak m(\log\Delta_n)\}.
\end{equation}
The expectations defining these centers use the unconditional laws, so
\(\widehat a_n\) is deterministic.

\begin{lemma}[Comparison of spatial normalizations]
\label{lem:space-normalization}
Under the metric-measure convergence assumption of the theorem,
\(\widehat a_n/a_n\to1\).
\end{lemma}

\begin{proof}
For every finite real \(Y\), bounded convergence shows that
\(F_Y(r)=\EE[\arctan(Y-r)]\) is continuous, strictly decreasing, and
has limits \(\pi/2\) and \(-\pi/2\) at the two ends of the real line.
Its zero therefore exists and is unique.  Directly,
\(\mathfrak m(Y+c)=\mathfrak m(Y)+c\).  If \(Y_j\) converges in law
to \(Y\), then \(F_{Y_j}(r)\to F_Y(r)\) for each fixed \(r\).
For each \(\varepsilon>0\), the values at
\(\mathfrak m(Y)-\varepsilon\) and
\(\mathfrak m(Y)+\varepsilon\) have opposite strict signs.  The same
eventually holds for \(F_{Y_j}\), trapping its zero between these two
points.  Thus \(\mathfrak m\) is continuous under convergence in law.

Diameter is continuous in the Gromov--Hausdorff topology.  Since
\(a_n\to0\), the assumed convergence gives
\[
 a_n\Delta_n\ \Longrightarrow\ \Delta,
 \qquad 0<\Delta<\infty\quad\hbox{almost surely}.
\]
Applying the preceding continuity statement after taking logarithms and
using the shift identity gives
\[
 \log\frac{a_n}{\widehat a_n}
 =\mathfrak m\bigl(\log(a_n\Delta_n)\bigr)
      -\mathfrak m(\log\Delta)\longrightarrow0.
\]
\end{proof}

\subsection{The geometric and analytic inputs}

We use \(\Sph\) for the coordinate sphere and \(\rho\) for its round
distance.  For each occurrence of a vertex, an incident edge, and an
incident face, form an equilateral triangle with vertices labelled by
these three objects.  Repeated incidences remain distinct.  Gluing the
\(4n\) triangles according to the cyclic incidences gives a topological
sphere with its piecewise Euclidean conformal structure.  Uniformize this
surface onto \(\Sph\), with the normalization of \cite{SC}, and write
\(z_n:V_n\to\Sph\) for the primal vertex locations.  The images of the
triangles are the embedded \emph{flags}; the primal edges are embedded
along their vertex-to-edge-midpoint sides.  This common coordinate will
retain both the intrinsic metric and the conformal energy.

\begin{proposition}[Inputs from the spectral companion]
\label{prop:inputs}
From every subsequence of sizes tending to infinity there is a further
subsequence and a representation coupling on which the following
properties hold almost surely, simultaneously.  The limiting sphere in
these coordinates has the ordinary unit-area \(\sqrt3\)-LQG law, with
area \(\mu_h\) and metric \(D_h\); the latter induces the round topology.

\begin{enumerate}[label=\textup{(\roman*)}]
\item \emph{Geometry in one coordinate.}  The vertex measures satisfy
\((z_n)_*\mu_n\to\mu_h\) weakly, and
\begin{align}
 \delta_n&:=\sup_{u,v\in V_n}
  \abs{a_nd_n(u,v)-D_h(z_n(u),z_n(v))}\longrightarrow0,
       \label{eq:vertex-distortion}\\
 \varepsilon_n&:=d_{\mathrm H}^{D_h}(z_n(V_n),\Sph)
       \longrightarrow0.\label{eq:vertex-density}
\end{align}
The maximum round diameter of the embedded flags tends to zero.
Consequently the maximum round diameter of an embedded primal edge tends
to zero as well.

\item \emph{One Hilbert space.}  There is an order
\(v_n(i)\), \(0\le i<2n\), of the primal corner occurrences and a
continuous tour \(\eta:[0,1]\to\Sph\) such that
\begin{equation}\label{eq:tour-tracking}
 \eta_*\dd t=\mu_h,\qquad
 \sup_{0\le t<1}
 \rho\bigl(z_n(v_n(\lfloor2nt\rfloor)),\eta(t)\bigr)
 \longrightarrow0.
\end{equation}
Each vertex \(v\) occurs \(\deg(v)\) times in the corner list.  For
\(\mu_h\)-almost every \(z\), the fiber \(\eta^{-1}(\{z\})\)
contains exactly one time.  With \(H=L^2([0,1],\dd t)\), the maps
\begin{align*}
 J_n:L^2(\mu_n)&\longrightarrow H,&
 (J_nu)(t)&=u(v_n(\lfloor2nt\rfloor)),\quad 0\le t<1,\\
 J:L^2(\mu_h)&\longrightarrow H,& Jf&=f\circ\eta
\end{align*}
are isometries, and \(J\) is unitary.  In particular,
\begin{equation}\label{eq:continuous-pullbacks}
 \norm{J_n(f\circ z_n)-Jf}_{L^\infty([0,1])}\longrightarrow0
 \qquad(f\in C(\Sph)).
\end{equation}

\item \emph{Centered inverse convergence.}  The closed continuum form
\(\E_h\) is a regular Dirichlet form and is the form of Liouville
Brownian motion with the stated half-energy clock.  Its operator \(G\)
has kernel exactly the constants.  For a nonnegative operator whose
kernel is the constants, write \(A^\dagger\) for its inverse on the
mean-zero subspace, extended by zero on constants.  Let \(\Pi\) be the
orthogonal projection onto the constants in \(H\), and put
\begin{equation}\label{eq:centered-inverses}
 K_n=J_n(2G_n)^\dagger J_n^*,\qquad
 K=J(2G)^\dagger J^*.
\end{equation}
Thus \(K_n\) is also zero on \(J_n(L^2(\mu_n))^\perp\).  These are
bounded positive self-adjoint operators, and
\begin{equation}\label{eq:inverse-convergence}
 \norm{K_n-K}_{\mathrm{HS}}\longrightarrow0.
\end{equation}

\item \emph{Local recovery with unit conductivity.}  There are ordinary
conformal chart domains whose interiors cover \(\Sph\).  Let
\(V\Subset W\Subset U\) be open sets with \(U\) compactly contained
in one of these charts, and let \(\chi\) be a real smooth function on
\(\Sph\) satisfying \(0\le\chi\le1\) and
\(\operatorname{supp}\chi\subset V\).  There are local vertex
functions \(\widetilde w_n\), defined on graphs containing all
incidences needed over \(W\), that converge uniformly to \(\chi\) on
a neighborhood of \(\overline W\), in the vertex locations, and satisfy
\begin{equation}\label{eq:local-recovery}
 \limsup_n\sum_{e:\,z_n(e)\in W}
       (d_e\widetilde w_n)^2
 \le\int_U\abs{\nabla\chi}^2\,\dd\mathrm{vol}.
\end{equation}
Here \(z_n(e)\) is any chosen point on the embedded edge, and the
integral is the unhalved conformal Dirichlet integral.  Every edge is
counted with its multiplicity.  If one endpoint lies in a fixed compact
subset of \(W\), then eventually its entire edge, both endpoints, and
all the required incidences lie in \(W\), uniformly over such edges.
\end{enumerate}
The assertions in \textup{(iv)} hold simultaneously for these buffered
tests on the coupled outcome, so the tests may be chosen after fixing
that outcome.
\end{proposition}

\begin{proof}[Source of the inputs and comparison of conventions]
The substantial assertions are imported from \cite{SC}.  Its
Proposition~2.4 supplies the common coordinate, the full flag mesh, and
the tours; its Theorem~5.2 supplies the local recovery assertion; its
Proposition~8.3 identifies the closed continuum form and centered inverse;
and its Lemma~9.1 and Equation~(9.11) give the Hilbert-space identification
and Hilbert--Schmidt convergence.  In the common extraction of its
Section~9, retain also the countable family of local constrained-energy
infima from its Lemma~5.1.  These are infima of finite sums of squared
edge differences under rational interval constraints on vertex values,
viewed in the compact space \([0,\infty]\).  Their joint limits supply
the simultaneous recovery statement of Theorem~5.2, and persist on
further subsequences.  Proposition~2.4 preserves the geometric
conclusions under this joint extraction, so all four assertions hold on
one coupling.

The model conventions agree exactly: the companion uses the same rooted
FK weight \(2^{\ell/2}\), counts all \(n\) primal edges, and pushes
normalized corner counting to \(\deg(v)/(2n)\).  Its all-edge network
therefore has operator \(B_n=2G_n\) by
Lemma~\ref{lem:graph-energy}.  Its limiting surface is the ordinary
unit-area sphere with the original marks forgotten.  Auxiliary marks and
rerootings used in the coupling change neither of these marginals.

For completeness, the companion first gives Hausdorff convergence of
the distance graphs
\[
 \{(z_n(u),z_n(v),\widehat a_nd_n(u,v)):u,v\in V_n\}
 \longrightarrow
 \{(x,y,D_h(x,y)):x,y\in\Sph\}
\]
in the product of the round topology and the ordinary topology on
\([0,\infty)\).  Uniform continuity of \(D_h\) on the compact
\(\Sph\times\Sph\) turns this into vanishing vertex distortion.
Approximating the triples \((x,x,0)\) in the opposite direction gives
round density of the vertex locations, hence \(D_h\)-density.
Lemma~\ref{lem:space-normalization} permits replacement of
\(\widehat a_n\) by \(a_n\): along this coupling the
\(\widehat a_n\)-scaled diameters are bounded, and the ratio of the
two scales tends to one.

The corner multiplicities give the isometry of \(J_n\), and the area
pushforward gives that of \(J\).  The almost-everywhere singleton-fiber
property makes \(J\) onto, as detailed in \cite[Lemma~9.1]{SC}.
Uniform tour tracking and uniform continuity of \(f\) give
\eqref{eq:continuous-pullbacks}.  The centered inverses in the companion
are those of the unhalved operators, which yields exactly
\eqref{eq:centered-inverses}.

Finally, the chart families in \cite[Proposition~2.4(iv)]{SC} use two
distinct exploration roots and separate auxiliary mark pairs.  Their
interiors cover the entire sphere, not merely the complement of one root.
The finite-sphere clause of its Theorem~5.2 gives
\eqref{eq:local-recovery} in the common coordinate, simultaneously for
continuous finite-energy inputs and buffered domains.  We have stated
only the smooth supported case needed here.  The vanishing flag mesh and
the strict buffers ensure that edges incident at the compact support are
indeed among the edges counted in \(W\).  Conformal and anticonformal
coordinate changes preserve the Dirichlet integral, and hence the
diffusion and its clock.
\end{proof}

Fix henceforth one outcome in the probability-one event of
Proposition~\ref{prop:inputs}, and relabel its further subsequence by
\(n\).  Until we return to random environments, all graphs, coordinates,
and the limiting measured conformal sphere are fixed.  The symbols
\(\PP\) and \(\EE\) will refer to the law and expectation of walks
sampled with fresh randomness on these fixed graphs, started from
\(\mu_n\), or of the limiting motion started from \(\mu_h\).

\section{Operators and finite-dimensional laws}\label{sec:operators}

Fix a realization in the probability-one event of
Proposition~\ref{prop:inputs}, and retain its extracted sequence. All limits in
this section and the next are on these fixed measured graphs and this fixed
measured sphere. Let $Y^n$ be the stationary chain with generator $-G_n$, and set
$\xi^n_t=z_n(Y^n_t)$. We first identify its finite-dimensional limits. The
centered inverses contain the required positive-time information, provided that
the constants and the orthogonal complements of the discrete spaces are
treated separately.

\begin{proposition}\label{prop:operators}
For every $s>0$ and $\alpha>0$, the following limits hold in operator norm on
$H$:
\begin{align}
 J_n e^{-sG_n}J_n^*&\longrightarrow J e^{-sG}J^*,
 \label{eq:semigroup-convergence}\\
 J_n\alpha(\alpha+G_n)^{-1}J_n^*
 &\longrightarrow J\alpha(\alpha+G)^{-1}J^*.
 \label{eq:resolvent-convergence}
\end{align}
Consequently the finite-dimensional laws of $\xi^n$ converge to those of
the stationary Liouville Brownian motion $B^h$.
\end{proposition}

\begin{proof}
Write $\Pi$ for the orthogonal projection onto the constants in $H$. On the
mean-zero part of $J_nL^2(\mu_n)$, the operator $2K_n$ is the inverse of $G_n$.
On the constants and on $(J_nL^2(\mu_n))^\perp$, it is zero. Therefore
\begin{equation}\label{eq:functional-calculus}
 J_n e^{-sG_n}J_n^*=\Pi+\psi_s(2K_n),
 \qquad
 \psi_s(y)=
 \begin{cases}e^{-s/y},&y>0,\\0,&y=0.\end{cases}
\end{equation}
In particular, both sides vanish on the orthogonal complement of the discrete
space, while both act as the identity on constants. The same identity holds
for $G$ and $K$, with no extra orthogonal complement since $J$ is unitary.
Hilbert--Schmidt convergence of $K_n$ implies operator-norm convergence and
places all their spectra in a common bounded interval. Since $\psi_s$ is
continuous on that interval, continuous functional calculus proves
\eqref{eq:semigroup-convergence}. Similarly,
\[
 J_n\alpha(\alpha+G_n)^{-1}J_n^*
   =\Pi+r_\alpha(2K_n),
 \qquad r_\alpha(y)=\frac{\alpha y}{1+\alpha y},
\]
and the corresponding identity for $G$ proves
\eqref{eq:resolvent-convergence}. Neither argument asserts operator-norm
convergence at time zero.

To identify finite-dimensional laws, take continuous real functions
$f_0,\ldots,f_k$ on $\Sph$ and times
$0\le t_0<t_1<\cdots<t_k$. Put $g_{n,i}=f_i\circ z_n$ and let $M_a$ denote
multiplication by a bounded function $a$ on $H$. By stationarity and the
Markov property,
\begin{equation}\label{eq:fdd-products}
 \EE\prod_{i=0}^k f_i(\xi^n_{t_i})
 =\big\langle 1,
 M_{J_ng_{n,0}}S_n(t_1-t_0)M_{J_ng_{n,1}}
 \cdots S_n(t_k-t_{k-1})M_{J_ng_{n,k}}1\big\rangle_H,
\end{equation}
where $S_n(s)=J_ne^{-sG_n}J_n^*$. Indeed, multiplication by $J_ng_{n,i}$
preserves the range of $J_n$ and represents multiplication by $g_{n,i}$ there.
The tour convergence in Proposition~\ref{prop:inputs} gives
$J_ng_{n,i}\to Jf_i$ uniformly, hence convergence of these multiplication
operators in norm. Equation~\eqref{eq:semigroup-convergence} now passes every
factor in \eqref{eq:fdd-products} to the limit. The resulting expression is
the same product expectation for $B^h$: its semigroup is $e^{-sG}$ and its
initial law $\mu_h$ is stationary because $G1=0$. Products of continuous
functions determine the laws on each finite product of the compact sphere.
Repeated observation times can be combined into a single test function, so
all finite-dimensional distributions converge.
\end{proof}

We shall also need the operator applied to a resolvent approximation. If
$f\in C(\Sph)$, $g_n=f\circ z_n$, and
\begin{equation}\label{eq:resolvent-functions}
 u_{n,\alpha}=\alpha(\alpha+G_n)^{-1}g_n,
 \qquad u_\alpha=\alpha(\alpha+G)^{-1}f,
\end{equation}
then, for every fixed $\alpha>0$,
\begin{equation}\label{eq:resolvent-graph-convergence}
 J_nu_{n,\alpha}\longrightarrow Ju_\alpha,
 \qquad J_nG_nu_{n,\alpha}\longrightarrow JGu_\alpha
 \quad\hbox{in }H.
\end{equation}
The first assertion follows from \eqref{eq:resolvent-convergence} and
$J_ng_n\to Jf$. The second follows from the resolvent equations
$G_nu_{n,\alpha}=\alpha(g_n-u_{n,\alpha})$ and
$Gu_\alpha=\alpha(f-u_\alpha)$.

\section{From energy recovery to path tightness}\label{sec:tightness}

Finite-dimensional convergence does not control excursions between observation
times. We obtain that control from the local energy recovery in
Proposition~\ref{prop:inputs}. First we extend buffered cutoff recoveries to
the whole graph. A reversible maximal inequality then makes their resolvent
approximations accurate along the entire stationary path. The bounded
generators of those approximations will control motion after stopping times.

\subsection{Uniform recovery and a stationary maximal inequality}

\begin{lemma}[Uniform cutoff recovery]\label{lem:uniform-recovery}
Let $\chi:\Sph\to[0,1]$ be smooth, and put $f=1-\chi$. Suppose
$\operatorname{supp}\chi\subset V\Subset W\Subset U$ for buffered conformal
chart regions to which the local recovery statement of
Proposition~\ref{prop:inputs} applies. There are real vertex functions $w_n$
such that
\begin{equation}\label{eq:uniform-recovery}
 \norm{w_n-f\circ z_n}_{\infty,V_n}\longrightarrow0,
 \qquad \limsup_n\E_n(w_n)\le\E_h(f).
\end{equation}
\end{lemma}

\begin{proof}
Apply local recovery to $\chi$ on $U$, obtaining $\widetilde\chi_n$ uniformly
close to $\chi\circ z_n$ at vertices in $W$. Choose $e_n\downarrow0$
bounding these uniform errors; decreasingness can be achieved by taking tail
suprema. Define
\[
 \chi_n(v)=
 \begin{cases}
 \min\{1,\max\{0,\widetilde\chi_n(v)-e_n\}\},&z_n(v)\in W,\\
 0,&z_n(v)\notin W,
 \end{cases}
 \qquad w_n=1-\chi_n.
\]
The clipping map is $1$-Lipschitz, and $\chi$ vanishes outside $W$, so
$\norm{\chi_n-\chi\circ z_n}_\infty\le2e_n$. More importantly,
$\chi_n(v)>0$ implies $\chi(z_n(v))>0$. Thus an edge with a nonzero difference
of $\chi_n$ has an endpoint in $\operatorname{supp}\chi\Subset V$.
The vanishing edge mesh and the buffer ensure that every such edge is
eventually wholly in $W$, with both endpoints there. Its difference is then
bounded in absolute value by the corresponding difference of
$\widetilde\chi_n$. Local recovery therefore gives
\[
 \limsup_n\E_n(w_n)
 =\limsup_n\frac12\sum_e(d_e\chi_n)^2
 \le\frac12\int_U\abs{\nabla\chi}^2\dd\mathrm{vol}
 =\E_h(f).
\]
The last equality uses the support condition and conformal invariance of
the two-dimensional Dirichlet integral. This proves
\eqref{eq:uniform-recovery}, including the factor $1/2$.
\end{proof}

The following estimate is a finite-state analogue of the forward--backward
martingale argument of Lyons and Zheng \cite{LyonsZheng1988}.
We include its proof to handle jumps explicitly and to make the uniformity
in the chain clear.

\begin{lemma}[Stationary maximal inequality]\label{lem:stationary-maximal}
Let $Y$ be a continuous-time Markov chain on a finite state space, with
generator $-A$ and reversible probability measure $\pi$, started from $\pi$.
For every real function $v$ on that space and every $T>0$,
\begin{equation}\label{eq:stationary-maximal}
 \EE\sup_{0\le t\le T}\abs{v(Y_t)}^2
 \le 2\norm{v}_{L^2(\pi)}^2+40T\langle v,Av\rangle_{L^2(\pi)}.
\end{equation}
The constants do not depend on the state space or its jump rates.
\end{lemma}

\begin{proof}
Write $q(x,y)$ for the off-diagonal jump rates. The process
\[
 N_t=v(Y_t)-v(Y_0)+\int_0^t Av(Y_s)\dd s
\]
is a square-integrable martingale starting from zero. Its quadratic variation
and stationarity give
\begin{align*}
 \EE\abs{N_T}^2
 &=T\sum_x\pi(x)\sum_{y\ne x}q(x,y)(v(y)-v(x))^2\\
 &=2T\langle v,Av\rangle_{L^2(\pi)}.
\end{align*}
Over $[0,T]$, reverse the path using its right-continuous modification in
reversed time. Reversibility and stationarity give the same chain law, so its
corresponding martingale $N'$ satisfies the same second-moment identity.
At every fixed $t\in[0,T]$, almost surely there is no jump at $t$ or at either
deterministic endpoint. Subtracting the forward and backward formulas yields
\[
 2\bigl(v(Y_t)-v(Y_0)\bigr)=N_t-(N'_T-N'_{T-t}).
\]
Use this identity first at a countable dense set of times, together with
$T$, and then use right continuity. It follows that
\[
 \sup_{t\le T}\abs{v(Y_t)}
 \le \abs{v(Y_0)}+\frac12\sup_{t\le T}\abs{N_t}
                       +\sup_{t\le T}\abs{N'_t}.
\]
If the two martingale suprema are $a$ and $b$, the square of the right side is
at most $2\abs{v(Y_0)}^2+a^2+4b^2$. Doob's $L^2$ inequality bounds each of
$\EE a^2$ and $\EE b^2$ by
$8T\langle v,Av\rangle_{L^2(\pi)}$, proving
\eqref{eq:stationary-maximal}.
\end{proof}

\subsection{Resolvent approximation along the whole path}

\begin{lemma}\label{lem:resolvent-approximation}
For a cutoff $f=1-\chi$ as in Lemma~\ref{lem:uniform-recovery}, write
$g_n=f\circ z_n$ and define $u_{n,\alpha}$ and $u_\alpha$ by
\eqref{eq:resolvent-functions}. For every $T>0$,
\begin{equation}\label{eq:pathwise-resolvent-approximation}
 \lim_{\alpha\to\infty}\limsup_n
 \EE\sup_{0\le t\le T}
       \abs{u_{n,\alpha}(Y^n_t)-g_n(Y^n_t)}=0.
\end{equation}
For every $n$ and $\alpha>0$, moreover,
\begin{equation}\label{eq:bounded-generator}
 0\le u_{n,\alpha}\le1,
 \qquad \norm{G_nu_{n,\alpha}}_\infty\le\alpha.
\end{equation}
\end{lemma}

\begin{proof}
Take the recoveries $w_n$ from Lemma~\ref{lem:uniform-recovery}. Their uniform
approximation, together with the tour convergence, gives $J_nw_n\to Jf$ in
$H$. For fixed $\alpha$, Equation~\eqref{eq:resolvent-graph-convergence}
therefore implies
\begin{align}
 &\limsup_n\bigl(
   \norm{w_n-u_{n,\alpha}}_{L^2(\mu_n)}^2
       +\E_n(w_n-u_{n,\alpha})\bigr)\notag\\
 &\hspace{15mm}\le
   \norm{f-u_\alpha}_{L^2(\mu_h)}^2+\E_h(f-u_\alpha).
 \label{eq:form-norm-approximation}
\end{align}
To see the energy assertion explicitly, expand
\[
 \E_n(w_n-u_{n,\alpha})
 =\E_n(w_n)-2\langle w_n,G_nu_{n,\alpha}\rangle_{L^2(\mu_n)}
            +\langle u_{n,\alpha},G_nu_{n,\alpha}\rangle_{L^2(\mu_n)}.
\]
The last two terms converge by \eqref{eq:resolvent-graph-convergence}, and
the first has limsup at most $\E_h(f)$. Since $u_\alpha$ lies in the operator
domain of $G$, the resulting expression is the form expansion of
$\E_h(f-u_\alpha)$.

The right side of \eqref{eq:form-norm-approximation} tends to zero. Indeed,
if $\nu_f$ is the spectral measure of $G$ at $f$, then it equals
\[
 \int_{[0,\infty)}(1+\lambda)
       \left(\frac{\lambda}{\alpha+\lambda}\right)^2\dd\nu_f(\lambda).
\]
Here $f$ belongs to the form domain, so $1+\lambda$ is integrable against
$\nu_f$; dominated convergence applies. Lemma~\ref{lem:stationary-maximal},
applied with $A=G_n$ to $w_n-u_{n,\alpha}$, now gives
\[
 \lim_{\alpha\to\infty}\limsup_n
 \EE\sup_{t\le T}\abs{w_n(Y^n_t)-u_{n,\alpha}(Y^n_t)}^2=0.
\]
Cauchy--Schwarz and the uniform bound on $w_n-g_n$ yield
\eqref{eq:pathwise-resolvent-approximation}.

Finally, the Markov semigroup representation
\[
 u_{n,\alpha}=\int_0^\infty\alpha e^{-\alpha t}
                              e^{-tG_n}g_n\dd t
\]
shows that $0\le u_{n,\alpha}\le1$. The resolvent identity
$G_nu_{n,\alpha}=\alpha(g_n-u_{n,\alpha})$ proves the remaining bound in
\eqref{eq:bounded-generator}.
\end{proof}

We now have approximations accurate for all times of the stationary path,
not merely at a fixed time. Their generator bound makes that accuracy useful
even when the starting time of an increment is selected from the path.

\subsection{Stopping-time increments}

\begin{proposition}\label{prop:path-limit}
On every fixed realization supplied by Proposition~\ref{prop:inputs}, and for
every $T>0$,
\[
 (\xi^n_t)_{0\le t\le T}\ \Longrightarrow\
 (B^h_t)_{0\le t\le T}
 \quad\hbox{in }D([0,T],\Sph)
\]
for the Skorokhod $J_1$ topology associated with the round metric $\rho$.
Every subsequential path limit is continuous.
\end{proposition}

\begin{proof}
Fix $d>0$. Choose finitely many round disks $O_1,\ldots,O_m$ covering
$\Sph$ and smooth cutoffs $f_i=1-\chi_i$ to which
Lemma~\ref{lem:uniform-recovery} applies, such that
\begin{equation}\label{eq:separating-cutoffs}
 f_i=0\text{ on }O_i,
 \qquad x\in O_i,\ \rho(x,y)\ge d\ \Longrightarrow\ f_i(y)=1.
\end{equation}
For example, take sufficiently small inner disks and slightly larger
support disks, all compactly contained in chart regions and all with
diameter smaller than $d$. Compactness gives a finite subcover, and there is
room for the recovery buffers outside each support disk.

Let $\tau,\sigma$ be stopping times for $Y^n$, with
$0\le\tau\le\sigma\le T$ and $\sigma-\tau\le s$. Let $I$ be the first index
for which $\xi^n_\tau\in O_I$. This index is measurable with respect to the
information $\mathcal F_\tau$ at time $\tau$. Since $0\le f_i\le1$,
\eqref{eq:separating-cutoffs} implies
\[
 \PP\{\rho(\xi^n_\tau,\xi^n_\sigma)\ge d\}
 \le\EE\bigl[f_I(\xi^n_\sigma)-f_I(\xi^n_\tau)\bigr].
\]
Denote by $u^{(i)}_{n,\alpha}$ the resolvent approximation of $f_i\circ z_n$.
Optional sampling of its generator martingale gives
\begin{align*}
 &\EE\bigl[u^{(I)}_{n,\alpha}(Y^n_\sigma)
                -u^{(I)}_{n,\alpha}(Y^n_\tau)\bigr]\\
 &\qquad=-\sum_{i=1}^m\EE\left[
    \mathbf1_{\{I=i\}}\int_\tau^\sigma
                    G_nu^{(i)}_{n,\alpha}(Y^n_t)\dd t\right]
 \le\alpha s.
\end{align*}
The equality is valid because each event $\{I=i\}$ lies in
$\mathcal F_\tau$; no independence of $I$ and the path is required.
The last inequality uses \eqref{eq:bounded-generator} and the fact that the
events $\{I=i\}$ partition the sample space. Consequently
\begin{align}
 \PP\{\rho(\xi^n_\tau,\xi^n_\sigma)\ge d\}
 &\le\alpha s+2\sum_{i=1}^m
 \EE\sup_{0\le t\le T}
   \abs{f_i(\xi^n_t)-u^{(i)}_{n,\alpha}(Y^n_t)}.
 \label{eq:aldous-bound}
\end{align}
This bound is uniform over all the stated stopping-time pairs. First take
$\limsup_n$, then let $s\downarrow0$ with $\alpha$ fixed, and finally let
$\alpha\to\infty$. Lemma~\ref{lem:resolvent-approximation} gives
\[
 \lim_{s\downarrow0}\limsup_n
 \sup_{\substack{\tau\le\sigma\le T\\0\le\sigma-\tau\le s}}
 \PP\{\rho(\xi^n_\tau,\xi^n_\sigma)\ge d\}=0.
\]
This is the stopping-time increment condition in Aldous's tightness
criterion \cite{Aldous1978}. Compact containment is automatic on $\Sph$,
so the laws are tight in $D([0,T],\Sph)$ with its $J_1$ topology.

Every jump of $\xi^n$ joins the endpoints of an embedded map edge.
Proposition~\ref{prop:inputs} therefore makes its largest jump tend to zero.
A jump of a $J_1$ limit must be approximated by jumps of comparable size;
hence every subsequential limit is continuous. Evaluation at fixed times is
continuous at continuous paths, so Proposition~\ref{prop:operators}
identifies the finite-dimensional laws of every such limit with those of
$B^h$. These finite-dimensional laws determine the continuous path law.
Tightness and uniqueness of the subsequential limit prove the proposition.
\end{proof}

\section{Cable paths and the joint limit}\label{sec:paths}

Proposition~\ref{prop:path-limit} identifies the walk in conformal
coordinates and with exponential holding times.  We now place the original
cable graphs and the limiting metric sphere in one compact space, replace
the exponential clock by the deterministic clock, and restore the
randomness of the environment.  Until the final proof, the realization
supplied by Proposition~\ref{prop:inputs} remains fixed.

\subsection{A common metric space}

The following construction uses only the vertex-distance comparison.  In
particular, it does not require an extension of the conformal coordinate
map to the interiors of cable edges.

\Needspace{8\baselineskip}
\begin{lemma}[A compact ambient space]\label{lem:ambient-space}
Let $C_n$ be finite connected cable graphs, with vertex sets $V_n$ and
all edge lengths equal to $a_n>0$, where $a_n\to0$.  Write $d_{C_n}$ for
their path metrics.  Let $(S,d)$ be a compact metric space, and suppose
that maps $z_n:V_n\to S$ satisfy
\begin{align*}
 \delta_n&:=\sup_{u,v\in V_n}
  \abs{d_{C_n}(u,v)-d(z_n(u),z_n(v))}\longrightarrow0,\\
 \varepsilon_n&:=d_{\mathrm H}^{S}(z_n(V_n),S)\longrightarrow0.
\end{align*}
There is a compact metric space $(Z,d_Z)$ containing isometric copies of
$S$ and every $C_n$, such that, with $c_n=\delta_n/2+a_n$,
\begin{align}
 \sup_{v\in V_n}d_Z(v,z_n(v))&\le c_n,
 \label{eq:vertex-bridges}\\
 d_{\mathrm H}^{Z}(C_n,S)&\le
  \max\{a_n+c_n,\varepsilon_n+c_n\}.
 \label{eq:cable-hausdorff}
\end{align}
If $\nu_n$ is a probability measure on $V_n$ and
$(z_n)_*\nu_n\Rightarrow\nu$ on $S$, then
$d_{\mathrm P}^{Z}(\nu_n,\nu)\to0$.
\end{lemma}

\begin{proof}
Start with the disjoint union of $S$ and the $C_n$.  Retain the metric
within each piece, and allow a crossing from $v\in V_n$ to $z_n(v)\in S$
at cost $c_n$.  Define distance as the infimum of the costs of finite
chains made of such crossings and movements within individual pieces.

This construction does not shorten any of the original metrics.  Indeed,
an excursion from $S$ through $C_n$, entering at $z_n(u)$ and returning
at $z_n(v)$, costs at least
\[
 2c_n+d_{C_n}(u,v)\ge d(z_n(u),z_n(v)).
\]
Replacing all such excursions by movements in $S$ proves the assertion
for pairs of points of $S$.  For points of a fixed $C_n$, first remove
excursions through the other cable graphs in this way.  Each remaining
excursion out of $C_n$, from a vertex $u$ to a vertex $v$, costs at least
\[
 2c_n+d(z_n(u),z_n(v))\ge d_{C_n}(u,v),
\]
so it too can be removed.  This also proves the assertion for points in
edge interiors.  Distinct pieces remain separated: every chain leaving
$C_n$ pays at least the positive cost $c_n$.  Thus the chain construction
defines a metric, not merely a pseudometric.

Every point of $C_n$ is at distance at most $a_n$ from a vertex, even
when its edge is a loop.  It is consequently at distance at most
$a_n+c_n$ from $S$.  Conversely, every point of $S$ is at distance at
most $\varepsilon_n+c_n$ from $C_n$.  These observations prove
\eqref{eq:vertex-bridges}--\eqref{eq:cable-hausdorff} on the union.  The
union is totally bounded: its tail lies arbitrarily close to the compact
space $S$, and each of the finitely many remaining cable graphs is
compact.  Its completion is the required compact space $Z$.

The coupling $v\mapsto(v,z_n(v))$ of $\nu_n$ and $(z_n)_*\nu_n$ has
all its mass at pairs at distance at most $c_n$, whence
\[
 d_{\mathrm P}^{Z}(\nu_n,\nu)
 \le c_n+d_{\mathrm P}^{Z}((z_n)_*\nu_n,\nu)\longrightarrow0.
\]
For the last convergence, weak convergence on the compact subspace $S$
is also weak convergence in $Z$.
\end{proof}

Apply Lemma~\ref{lem:ambient-space} to
$C_n=(\overline M_n,a_nd_n)$ and $(S,d)=(\Sph,D_h)$, with
$\nu_n=\mu_n$.  The hypotheses are precisely the metric and measure
conclusions of Proposition~\ref{prop:inputs}; vertex distances in a
cable graph equal the corresponding graph distances.  We henceforth
regard these spaces as subsets of the resulting $Z$.

Because $\rho$ and $D_h$ induce the same topology on the compact sphere,
the identity map between the two metric spheres is uniformly continuous
in both directions.  Applying the same time changes in the definition
of the Skorokhod $J_1$ topology therefore transfers
Proposition~\ref{prop:path-limit} from $\rho$ to $D_h$.  Moreover,
\eqref{eq:vertex-bridges} gives the deterministic bound
\[
 \sup_{0\le t\le T}d_Z(Y^n_t,\xi^n_t)\le c_n.
\]
It follows that, for every $T<\infty$,
\begin{equation}\label{eq:ambient-path-limit}
 Y^n\ \Longrightarrow\ B^h
 \quad\text{in }D([0,T],Z)\text{ with the }J_1\text{ topology}.
\end{equation}

We shall use one consequence of the continuity of $B^h$.  For a path
$y:[0,T]\to Z$, set
\[
 \omega_T(y,s)=
 \sup_{\substack{t,t'\in[0,T]\\\abs{t-t'}\le s}}
    d_Z(y_t,y_{t'}).
\]
Then, for every $\eta>0$,
\begin{equation}\label{eq:continuous-limit-modulus}
 \lim_{s\downarrow0}\limsup_n
 \PP\bigl[\omega_T(Y^n,s)>\eta\bigr]=0.
\end{equation}
To see this directly, use the Skorokhod representation theorem
\cite[Theorem~6.7]{Billingsley1999} on the Polish space $D([0,T],Z)$ in
\eqref{eq:ambient-path-limit}.  Convergence in $J_1$ to a continuous
path is uniform: composing with the convergent time changes changes that
limiting path uniformly little; see also
\cite[Section~12, p.~124]{Billingsley1999}.  On this representation,
\[
 \omega_T(Y^n,s)\le
 2\sup_{t\le T}d_Z(Y^n_t,B^h_t)+\omega_T(B^h,s).
\]
Uniform continuity of each limiting path proves
\eqref{eq:continuous-limit-modulus}.  The resulting probability statement
depends only on the original path laws, not on the chosen representation.

\subsection{Removing the exponential clock}

The clock comparison must count every attempted step, including a loop
traversal that does not change the vertex of the continuous-time chain.

\begin{lemma}[The attempt clock]\label{lem:clock}
Fix a realization supplied by Proposition~\ref{prop:inputs}, for which
\eqref{eq:ambient-path-limit} holds.  Let $X^n$ be the stationary discrete
walk on $M_n$, including its selected half-edges, and let $\overline X^n$
be its cable interpolation.  Independently of the walk, let
$(e_{n,k})_{k\ge1}$ be
independent exponential random variables of mean $1/n$, and put
\[
 S_{n,0}=0,\qquad S_{n,k}=\sum_{j=1}^k e_{n,j},\qquad
 N_n(t)=\max\{k:S_{n,k}\le t\}.
\]
Then $Y^n_t=X^n_{N_n(t)}$ has generator $-G_n$ and
$Y^n_{S_{n,k}}=X^n_k$.  In the ambient space of
Lemma~\ref{lem:ambient-space}, for every $T<\infty$,
\begin{equation}\label{eq:clock-comparison}
 \sup_{0\le t\le T}
 d_Z(\overline X^n_{nt},Y^n_t)\longrightarrow0
 \quad\text{in probability}.
\end{equation}
\end{lemma}

\begin{proof}
Attempts occur at rate $n$ and apply the transition kernel $P_n$,
so the generator is $n(P_n-I)=-G_n$.  This remains true when an
attempt traverses a loop.  The identity at $S_{n,k}$ follows from the
right-continuous definition of $N_n$.

Fix $T$ and write
\[
 \Delta_{n,T}=
 \max_{0\le k\le\lfloor nT\rfloor}\abs{S_{n,k}-k/n}.
\]
The centered partial sums form a martingale.  Since
$\operatorname{Var}(e_{n,k})=n^{-2}$, the $L^2$ maximal inequality gives
\begin{equation}\label{eq:clock-deviation}
 \EE[\Delta_{n,T}^2]
 \le \frac{4\lfloor nT\rfloor}{n^2}\longrightarrow0.
\end{equation}
For $k=\lfloor nt\rfloor$, the point $\overline X^n_{nt}$ lies on the
cable edge traversed from $X^n_k$, at distance at most $a_n$ from that
vertex.  Thus
\[
 d_Z(\overline X^n_{nt},Y^n_t)
 \le a_n+d_Z(Y^n_{S_{n,k}},Y^n_t),
 \qquad
 \abs{S_{n,k}-t}\le\Delta_{n,T}+1/n.
\]
The spatial bound holds for loop edges as well as nonloop edges.
On the event $\Delta_{n,T}+1/n\le s\le1$, all the times appearing here
lie in $[0,T+1]$.  Consequently, for $a_n<\eta/2$,
\[
 \PP\left[\sup_{t\le T}
 d_Z(\overline X^n_{nt},Y^n_t)>\eta\right]
 \le \PP[\Delta_{n,T}+1/n>s]
   +\PP[\omega_{T+1}(Y^n,s)>\eta/2].
\]
First take the upper limit in $n$, using
\eqref{eq:clock-deviation}, and then let $s\downarrow0$, using
\eqref{eq:continuous-limit-modulus} on $[0,T+1]$.  This proves
\eqref{eq:clock-comparison}.
\end{proof}

Combining Lemma~\ref{lem:clock} and \eqref{eq:ambient-path-limit} gives
convergence of the cable paths to $B^h$ in $D([0,T],Z)$.  Both the cable
paths and the limiting paths are continuous.  The continuous path space
is closed in the $J_1$ topology: a convergent sequence of continuous
paths, after the defining time changes, converges uniformly and hence
has a continuous limit.  Its induced topology is the uniform topology,
by the same time-change argument used above.  Thus in fact
\begin{equation}\label{eq:cable-finite-horizon}
 (\overline X^n_{nt})_{0\le t\le T}
 \ \Longrightarrow\ (B^h_t)_{0\le t\le T}
 \quad\text{in }C([0,T],Z).
\end{equation}

This convergence holds on every integer time horizon, and therefore also
in $C([0,\infty),Z)$ with the topology of uniform convergence on compact
time intervals.  Here are the compactness details.  For each integer
$m\ge1$ and $\epsilon>0$, finite-horizon tightness supplies a compact
set $K_m\subset C([0,m],Z)$ such that the probability of leaving $K_m$
is at most $\epsilon2^{-m}$, uniformly in $n$.  The set of paths whose
restriction belongs to $K_m$ for every $m$ is compact in
$C([0,\infty),Z)$, by a diagonal subsequence argument, and has
probability at least $1-\epsilon$.  Every subsequential limit has, by
the continuous restriction maps, the law of $B^h$ on every $[0,m]$.
These restrictions determine the entire path law.  We have proved
\begin{equation}\label{eq:cable-local-limit}
 (\overline X^n_{nt})_{t\ge0}
 \ \Longrightarrow\ (B^h_t)_{t\ge0}
 \quad\text{in }C([0,\infty),Z).
\end{equation}

\subsection{Restoring the random environment}

We have obtained the conditional path limit for every fixed realization
in the probability-one event of Proposition~\ref{prop:inputs}.  We now
combine it with convergence of the underlying measured spaces.

\begin{proof}[Proof of Theorem~\ref{thm:main}]\label{proofmain}
Start from an arbitrary subsequence and take the coupled further
subsequence of Proposition~\ref{prop:inputs}.  Fix a realization in
its probability-one event and a finite integer $r\ge1$.
Conditional on this realization, the $r$ cable paths are independent
with the common law in \eqref{eq:cable-local-limit}.  Their joint law
therefore converges to that of $r$ independent copies of $B^h$ in
$C([0,\infty),Z)^r$.  For example, a Skorokhod representation of
\eqref{eq:cable-local-limit}, repeated on $r$ independent probability
spaces, couples the product laws with almost sure convergence of every
coordinate.

Lemma~\ref{lem:ambient-space} also gives Hausdorff convergence of the
spaces and Prokhorov convergence of their measures in the same $Z$.
Under the preceding product coupling, the curves converge uniformly on
every compact time interval.  By the definition of the curve-decorated
Gromov--Hausdorff--Prokhorov topology, the decorated spaces consequently
converge almost surely under that coupling.  In particular, for any
bounded continuous function $F$ on the space of such decorated objects,
the expectation of $F$ over the $r$ walks on this fixed environment
converges to its expectation over $r$ independent stationary Liouville
Brownian motions on the fixed limiting sphere.

These expectations are the conditional expectations obtained by sampling
walks with fresh randomness after the environment coupling.  The walk
kernel on a finite graph depends only on the map, not on the auxiliary
coordinates, marks, or tours used in that coupling; the limiting kernel
is the diffusion law specified by the quantum sphere.  We have proved
pointwise convergence of these conditional expectations on the
probability-one event.  The compact space $Z$, the path couplings, and
the recovery functions used earlier are witnesses for this pointwise
statement.  They need not be chosen measurably as the environment varies.
The conditional expectations from these sampling kernels are measurable,
and bounded convergence now gives convergence of the unconditional
expectations of $F$ along the coupled further subsequence.

Thus every subsequence admits a further subsequence with the asserted
joint limit.  This proves convergence along the full sequence, for
each fixed $r$.  Lemma~\ref{lem:clock} used the deterministic
acceleration $b_n=n$, which has exactly the time normalization in the
theorem.
\end{proof}

\bibliographystyle{amsplain}
\begingroup\raggedright
\bibliography{references}
\endgroup
\end{document}
